\documentclass{article}
\usepackage[margin=1in]{geometry}
\usepackage{amsmath,amssymb,amsthm,booktabs,array}
\usepackage[T1]{fontenc}
\usepackage{xurl}
\usepackage[hidelinks]{hyperref}
\setlength{\emergencystretch}{2em}
\newtheorem{proposition}{Proposition}
\newcommand{\id}{\mathbb I}
\newcommand{\tr}{\operatorname{Tr}}
\newcommand{\gap}{\operatorname{gap}}
\title{Spectral-gap literature audit:\\open questions and tractable targets}
\author{TNLean research audit}
\date{16 September 2026}
\begin{document}
\maketitle

\section{Decision}

The strongest finding is a negative answer to the unrestricted
\emph{value-convergence} question for the literal truncated,
locally translation-invariant (LTI) spectral-gap hierarchy of Rai, Kull,
Emonts, Tura, Schuch, and Baccari~\cite{LTI}.
For the commuting ferromagnetic Ising interaction, every finite level has
value one, whereas the periodic-boundary (PBC) chain gap is two:
\begin{equation}
  \delta_{\mathrm{LTI}}(n)=1,
  \qquad \Delta_N^{\mathrm{PBC}}=2.
  \label{eq:verdict}
\end{equation}
Section~\ref{sec:obstruction} proves this directly from the paper's
certificate definition. The argument has been independently checked
against the primary source and by finite exact enumeration. It is an
internal audit deduction, not a Lean-checked result; prior publication
of this observation has not been established or excluded.

This does not contradict a proved theorem in the paper. The authors pose
convergence as an open question. It also does not disprove eventual
\emph{detection} of a positive gap: the example is certified positive at
once. The recommended next research step is to settle the provenance and
scope of this negative answer before developing a larger certificate
method. A new method would have to address the discarded products of
spatially separated excitations; increasing telescopic freedom alone
cannot remove this obstruction.

Two other results of the audit prevent redundant projects. First, eventual
positive certification of standard gapped MPS parents follows from existing
finite-size criteria under their open-boundary (OBC) chain hypotheses; it is not a new
solution of general hierarchy completeness. Second, the July 2026 paper
of Rozm\'an, Moln\'ar, and Schuch~\cite{RMS} already computes the relevant
MPS ground-space angle through a fixed-dimensional virtual-space matrix.
Efficient overlap computation and basic MPS-aware certification are
therefore not unoccupied research directions.

\section{Verified sources and search scope}
\label{sec:sources}

The version-pinned primary texts used for mathematical claims are:
\begin{enumerate}
  \item \textbf{Rai et al.}~\cite{LTI}, arXiv:2411.03680v3,
  1 April 2026; published in \emph{Quantum} \textbf{10}, 2065,
  13 April 2026. The arXiv history lists v1 on 6 November 2024 and v2
  on 31 March 2026. All six authors, including Schuch, appear in the v1
  PDF as well as v3. This is not a later author addition.
  \item \textbf{Rozm\'an, Moln\'ar, and Schuch}~\cite{RMS},
  \emph{Lower Bounds on Spectral Gaps of Parent Hamiltonians via Tensor
  Networks}, arXiv:2607.19078v1, 21 July 2026. The identifier, title,
  authors, and version were independently verified.
  \item \textbf{Rao}~\cite{Rao}, arXiv:2510.08427. Its convergent
  hierarchy concerns a fixed finite qubit system, not thermodynamic
  completeness of the original LTI hierarchy. Its basic energy-gap
  construction uses the two lowest eigenvalues counting multiplicity,
  which must not be confused with the smallest positive eigenvalue
  above a degenerate ground space.
  \item \textbf{Xu et al.}~\cite{Xu}, arXiv:2606.03836v2,
  revised 31 July 2026, establish convergent certified \emph{upper}
  bounds on a bulk gap. Their introduction still describes completeness
  of the frustration-free lower-bound hierarchy as open. Their
  Supplementary Information, Definition 1.6, Remark 1.8, Proposition 1.9,
  and Section 4.1, distinguish local GNS non-degeneracy from global
  symmetry-breaking degeneracy. This is not a solution of LTI
  lower-bound completeness.
  \item \textbf{Lemm--Mozgunov}~\cite{LM} and
  \textbf{Warzel--Young}~\cite{WY} provide boundary-sensitive
  finite-size criteria and examples distinguishing boundary and bulk
  energy scales. These are relevant to interpreting, not resolving,
  the LTI completeness question.
\end{enumerate}

The search followed the Quantum cited-by list, relevant arXiv records,
version-pinned primary texts, and references in the July paper. This is
an audit as of 16 September 2026, not an exhaustive bibliometric search.
The cited-by service warns of incomplete coverage; some searches returned
no content and arXiv requests encountered rate limits. A failure to find a
later solution or a prior Ising observation is not evidence that neither
exists. Rendering dates on third-party HTML conversions were not treated
as version dates.

\paragraph{Targeted provenance check.}
The inspected texts of Refs.~\cite{LTI,RMS,Rao,Xu} did not reveal this
specific separated-Ising-wall obstruction. Public GitHub repository
search for \texttt{2411.03680} returned no repositories; issue search
returned a third-party survey, not an author discussion. No authoritative
original implementation was located in the inspected paper or publication
page. SciRate discussion was inaccessible. General-web queries combining
\texttt{2411.03680}, \texttt{Ising}, \texttt{convergence},
\texttt{LTI}, \texttt{domain walls}, and \texttt{counterexample}
returned no usable content and are not counted as successful zero-hit
searches. In particular, the public repository linked by Ref.~\cite{Xu}
is not evidence about the implementation of Ref.~\cite{LTI}.
The conclusion below concerns the published formulas, not uninspected code.

\section{The two distinct completeness questions}
\label{sec:questions}

Let the single-site Hilbert space be $\mathbb C^d$. A positive interaction
$h$ has range $k$ and translates $h_i$ on a periodic chain of $N$ sites.
Assume that
\begin{equation}
  H_N=\sum_{i=1}^N h_i
\end{equation}
is frustration free with ground energy zero. Write $\Delta_N$ for its
smallest strictly positive eigenvalue. These are the conventions in
Ref.~\cite{LTI}, Section 2; ground-state uniqueness is not assumed.

At an admissible support level $n$, using the sufficient condition
$n\geq2k+1$ stated in the general-range construction, the paper discards
sufficiently distant products from $H_N^2$. In the printed convention of its equations (4)--(9),
write the resulting operator and local generating term as
\begin{align}
  [H_N^2]_n
    &=\sum_i\left(h_i^2+
      \sum_{1\leq s<n-k}\{h_i,h_{i+s}\}\right),
      \label{eq:truncated}\\
  g_n(\delta)
    &=h_1^2+\sum_{1\leq s<n-k}\{h_1,h_{1+s}\}-\delta h_1,
      \label{eq:generator}
\end{align}
where $\{A,B\}=AB+BA$ and site indices in the first line are cyclic.
The local terms in $g_n$ are embedded into $n$ sites with identities on
unused sites. The level value, Ref.~\cite{LTI}, equation (11), is
\begin{equation}
  \delta_{\mathrm{LTI}}(n)
  =\sup\left\{\delta:\ \exists Y=Y^\dagger,\quad
    g_n(\delta)+\id\otimes Y-Y\otimes\id\succeq0\right\},
  \label{eq:lti}
\end{equation}
where $Y$ acts on $n-1$ sites. Supremum notation avoids an unnecessary
attainment assumption; the example below attains its optimum.
For sufficiently large rings, in particular $N\geq2n$ in the paper's
stated regime, summing a feasible local certificate gives
\begin{equation}
  [H_N^2]_n-\delta H_N\succeq0.
  \label{eq:global-certificate}
\end{equation}
The omitted products are positive, giving $H_N^2\succeq\delta H_N$.

The discussion in Ref.~\cite{LTI}, Section 5, asks whether the hierarchy
\begin{quote}
``always detects a gap at some finite level''
\end{quote}
and whether its certified gap
\begin{quote}
``always converges to the true gap''.
\end{quote}
The second question also appears after equation (37) in Section 3.2.
These are different assertions:
\begin{align}
  \inf_{N\geq N_0}\Delta_N>0
    &\ \Longrightarrow\ \exists n:\delta_{\mathrm{LTI}}(n)>0,
    \label{eq:detection}\\
  \lim_{n\to\infty}\delta_{\mathrm{LTI}}(n)
    &\ \stackrel{?}{=}\lim_{N\to\infty}\Delta_N,
    \label{eq:value-convergence}
\end{align}
with existence of the right-hand limit required in
\eqref{eq:value-convergence}. A finite-volume statement must not be silently
replaced by a statement about an unspecified infinite-volume GNS
representation.

\section{A separated-domain-wall obstruction}
\label{sec:obstruction}

Let $Z=\operatorname{diag}(1,-1)$ be the Pauli matrix and set
\begin{equation}
  h_i=\frac{\id-Z_iZ_{i+1}}2.
  \label{eq:ising}
\end{equation}
These are commuting positive projections on qubits. For a binary
configuration $x=(x_1,\ldots,x_N)$, define its wall indicators by
\begin{equation}
  w_i(x)=\boldsymbol{1}_{\{x_i\neq x_{i+1}\}},
  \qquad x_{N+1}=x_1.
\end{equation}
Then $H_N$ is diagonal and its eigenvalue on $|x\rangle$ is
$\sum_iw_i(x)$.

\begin{proposition}
For every $N\geq3$, the periodic Hamiltonian with interaction
\eqref{eq:ising} has gap $\Delta_N=2$. For every admissible finite
level of \eqref{eq:lti}, its hierarchy value is exactly one.
\end{proposition}
\begin{proof}
Every binary configuration on a ring has an even number of walls.
There are two zero-wall configurations, and a nonconstant interval of
ones in a background of zeros has exactly two walls. Consequently
\begin{equation}
  \ker H_N=\operatorname{span}\{|0^N\rangle,|1^N\rangle\},
  \qquad \Delta_N=2.
\end{equation}
Since $h_1^2=h_1$ and products of commuting positive projections are
positive,
\begin{equation}
  g_n(1)=2\sum_{1\leq s<n-k}h_1h_{1+s}\succeq0.
\end{equation}
Thus $\delta=1$, $Y=0$ is feasible.

Conversely, fix $n$ and take $N=4n$ with configuration
$x=0^{2n}1^{2n}$. Its two walls are separated by $2n$ in each cyclic
direction, so none of the retained products contains both walls.
Therefore
\begin{equation}
  H_N|x\rangle=2|x\rangle,
  \qquad [H_N^2]_n|x\rangle=2|x\rangle.
\end{equation}
For any feasible $\delta,Y$, the telescopic terms cancel as operators
when translated around the ring. Equation~\eqref{eq:global-certificate}
then implies
\begin{equation}
  0\leq\langle x|([H_N^2]_n-\delta H_N)|x\rangle
  =2-2\delta.
\end{equation}
Hence $\delta\leq1$, including for non-diagonal $Y$, and the feasible
value one is optimal.
\end{proof}

The order of quantifiers matters: each fixed level must certify every
sufficiently large ring, so a ring chosen as a function of that level is
a valid obstruction. No interchange of limits is used. The example
answers \eqref{eq:value-convergence} negatively for the literal hierarchy,
but satisfies \eqref{eq:detection}. It is a two-sector, GHZ-type parent,
not a single normal/injective-MPS parent.

There is a cutoff discrepancy between the paper's general generating-term
formula and its nearest-neighbor dual formula: the last included neighbor
differs by one. The proof is unchanged under either convention, because
the two walls are farther apart than the full support. Taking $n\geq5$
also satisfies the conservative general condition $n\geq2k+1$ with $k=2$.

\paragraph{Verification status.}
Two independent source inspections checked the assumptions, periodic-size
quantifier, and degeneracy convention. A separate exact enumeration of all
binary configurations gave periodic gap two and minimum ratio
$\langle[H_N^2]_n\rangle/\langle H_N\rangle=1$ for
$(n,N)=(5,12)$ and $(6,14)$, among other sizes. An independent check also
constructed feasible local dual marginals. These computations support,
but do not replace, the all-level proof above. No Lean proof or full
library build was performed.

\section{What the newer literature already accomplishes}
\label{sec:occupied}

\paragraph{Known finite-size inclusion.}
For nearest-neighbor projector interactions, Ref.~\cite{LTI}, equations
(23), (34), and (36), and Appendix B.1, give same-level Knabe inclusion
in the nearest-neighbor cutoff convention used there. With the strict
literal cutoff in \eqref{eq:truncated}, the safe statement is shifted:
\begin{equation}
  \delta_{\mathrm{LTI}}(n+1)\geq
  \frac{n-1}{n-2}\left(\epsilon_n-\frac1{n-1}\right),
  \label{eq:knabe-inclusion}
\end{equation}
where $\epsilon_n$ is the open-chain gap on $n$ sites, with $n$ large
enough for the stated criteria. Indeed, the conventional $n$-site
certificate retains separations through $n-2$; tensoring it with one
identity gives exactly the strict-cutoff $(n+1)$-site certificate.
The corresponding telescopic variable is $Y\otimes\id$. A uniform positive
lower bound on $\epsilon_n$ therefore already proves eventual detection.
If open and periodic limiting gaps exist and coincide, the lower bound
\eqref{eq:knabe-inclusion} and periodic certificate validity also give
value convergence by squeezing. These are not new completeness targets.
The martingale inclusion in Appendix B.3 is stated for a coarse-grained
projector Hamiltonian; transfer back to the original interaction is a
separate operator comparison.

\paragraph{Ground-space-aware certificates are already present.}
Ref.~\cite{LTI}, Appendix D, equations (115)--(117), restricts the auxiliary
operator to $P_{n-1}^{\perp}YP_{n-1}^{\perp}$, where $P_{n-1}$ projects
onto the open-chain ground space. The paper explicitly permits constructing
this projector from the MPS ground-space representation. Such a restriction
may lower the SDP optimum. Neither this ansatz nor basic numerical
certificate repair should be proposed as new mathematics.

\paragraph{The July 2026 improvement.}
Ref.~\cite{RMS}, Lemmas 1 and 2 together with Theorems 3 and 4, obtains
\begin{equation}
  \Delta(H_N)\geq
  (1-2\gamma_{p,q}) \Delta(\widetilde h),
  \qquad \gamma_{p,q}=\sqrt{\lambda_2(\Xi)}.
  \label{eq:rms-bound}
\end{equation}
Here $p,q$ specify overlapping blocks, $\widetilde h$ is the weighted
blocked interaction on $r=2p+q$ sites, and $\Xi$ has virtual dimension
$D^4$ for MPS bond dimension $D$. The symbol $\lambda_2$ denotes the
largest eigenvalue distinct from the ground-space eigenvalue one.
The paper already removes exponential physical-size dependence from this
overlap calculation. Its remaining bottleneck is the local gap of
$\widetilde h$, a $d^r\times d^r$ operator. The arithmetic estimates in
Section 4 give full diagonalization cost $O(d^{3r})$; these estimates are
not, by themselves, bit-complexity bounds for exact certification.

Section 4.4 shows convergence of the resulting periodic lower bounds to
an \emph{open-boundary} limiting value,
$\limsup_m\Delta(H_m^{\mathrm{OBC}})$, under the specified asymmetric
blocking. It does not prove convergence to the periodic gap. Section 6
proposes extending reshuffling to ``a telescopic sum of arbitrary
operators.'' That suggestion is not a conjecture, and telescopic
optimization alone is already central to Ref.~\cite{LTI}.

\paragraph{Published numerical benchmarks.}
The following bounds must retain the interaction normalizations of their
respective sources:
\begin{center}
\begin{tabular}{p{0.28\linewidth}p{0.30\linewidth}p{0.29\linewidth}}
\toprule
Model and source & Reported bound & Resource or scope \\
\midrule
AKLT, Ref.~\cite{LTI}, Section 4.1
& $\Delta\geq0.34976$ & LTI support $n=6$ \\
AKLT, Ref.~\cite{RMS}, (33)--(34)
& PBC $0.2827$; OBC $0.2497$ & Best reported search through $r=14$ \\
$\mathrm{SU}(3)$ adjoint $\mathbf8$, Ref.~\cite{RMS}, (46)--(47)
& OBC $0.3809$; PBC $0.3917$ & Blocks through $r=6$ \\
$\mathrm{SU}(3)$ $\mathbf{10}$, Ref.~\cite{RMS}, (54)--(55)
& OBC $0.2158$; PBC $0.2204$ & Blocks through $r=6$ \\
\bottomrule
\end{tabular}
\end{center}
The two $\mathrm{SU}(3)$ gappedness questions discussed in the July paper
are solved there; they are not available conjectures. Improving the weaker
July AKLT bound alone is not competitive with the earlier LTI result.
Comparisons across methods also require the same boundary conditions,
Hamiltonian normalization, and a declared cost model, not merely a matching
support label.

\section{Existing formal results and the missing mathematics}
\label{sec:lean}

Source inspection was performed at TNLean commit \texttt{4129cbabe} and
its QICLean dependency. This is not a new compilation or an axiom audit.
Paths below are relative to \path{TNLean/MPS/ParentHamiltonian/}.
\begin{itemize}
  \item \path{PrimitiveGaugeExistence.lean}, lines 232--254:
  \path{exists_parentHamiltonianES_gap_of_isNormal} chooses an interaction
  range and proves a uniform finite-volume periodic gap for a normal
  tensor. It does not cover every prescribed range or general multiblock
  tensors, and is not an infinite-volume GNS theorem.
  \item \path{Martingale/OpenParentGap.lean}, lines 230--242:
  \path{IsPrimitiveMPS.exists_openParentHamiltonianES_gap} chooses a range
  $l+1$ and gives $\gamma=(1-\varepsilon\sqrt{l+1})^2$ for the open parent.
  \item \path{Martingale/FiniteRangeKnabeGap.lean}, lines 47--57:
  \path{parentHamiltonianES_gap_of_openParentHamiltonianES_gap} accepts a
  prescribed range $R$ with $1\leq R\leq m$, an open-window gap $\gamma$
  on $m+R-1$ sites, and $m\gamma>(R-1)^2$. It yields
  \begin{equation}
    \delta=\frac{m\gamma-(R-1)^2}{m-R+1}>0,
    \qquad N\geq2m.
  \end{equation}
  \item \path{Martingale/BlockedGap.lean}, lines 412--431, proves a gap
  $1/8$ for a suitably blocked range-two parent. The parent interaction
  changes; $1/8$ is not a universal improvement for the original Hamiltonian.
  \item \path{FNWGeometricDefect.lean}, lines 69--101, proves the conditional
  overlap estimate, for $c\lambda^l<1$ and an overlap length $l$ at
  least the supplied injectivity length,
  \begin{equation}
    \varepsilon_l\leq
    \frac{c\lambda^l(1+c\lambda^l)}{1-c\lambda^l}.
  \end{equation}
  The decay prefactor $c$ and chosen threshold length are supplied
  existentially. These are established FNW/Nachtergaele estimates,
  not stronger new constants.
\end{itemize}
QICLean supplies semidefinite weak duality, positive-operator gap transfer,
and finite-range Knabe inequalities. No implementation of the specific
LTI hierarchy or its completeness theorem was found in the inspected
sources. An exact formalization of Section~\ref{sec:obstruction} would
mostly require the hierarchy definition, telescopic summation, and binary
wall-counting, rather than more parent-Hamiltonian analytic estimates.

For a quantitative improvement to \eqref{eq:rms-bound}, the missing
mathematics is different: a stronger bound on the same local interaction,
or a controlled reduction of the one-block gap computation. The existence
of a small-bond-dimension ground state alone does not supply a complete
low-energy representation or justify discarding excited sectors.

\section{Ranked research decisions}
\label{sec:targets}

\paragraph{1. Preferred: resolve and contextualize the LTI convergence question.}
The exact paper target is Ref.~\cite{LTI}, Section 5, value convergence.
Section~\ref{sec:obstruction} gives a negative answer in its stated
periodic, possibly degenerate setting. The next finite deliverable is a
source-faithful, independently checkable statement of this obstruction,
with provenance checked against the authors' code and later discussion.
The current evidence supports the mathematics, not a claim of first
publication. If the observation is already known, record it and do not
repackage it as a new conjecture solution. A substantial extension would
characterize what the truncated hierarchy converges to or change the
certificate family to recover separated-excitation costs. Neither extension
is established by this audit.

\paragraph{2. Genuine open question, higher risk: eventual positivity.}
The remaining question \eqref{eq:detection} is explicitly posed in
Ref.~\cite{LTI}, Section 5. No resolution was found in the inspected
follow-ups. Normal-MPS cases with uniform open-chain gaps are not the hard
part, by \eqref{eq:knabe-inclusion}. A meaningful attack must address cases
not covered by that implication or produce a counterexample to detection.
The current TNLean capstones do not establish such a result. This is an
important target, but not presently justified as a short project.

\paragraph{3. Quantitative fallback: the original $\mathrm{SU}(3)$
$\mathbf{10}$ interaction.}
Fix the interaction of Ref.~\cite{RMS}, equation (52), its normalization,
and periodic boundaries. With generators $J^a$ in the ten-dimensional
representation normalized by $\sum_{a=1}^8(J^a)^2=6\id$, set
\begin{equation}
  X_i=\sum_{a=1}^8J_i^aJ_{i+1}^a,
  \qquad H_N^{\circ}=\sum_i(X_i^2+5X_i+6\id).
\end{equation}
Do not substitute the spectrally flattened projector interaction when
comparing gap values. The explicit benchmark is the bound $0.2204$ in
equation (55), obtained using blocks through six sites. A worthwhile
improvement must substantially exceed that bound under a comparable cost
model, or recover comparable quality at substantially lower cost. A
slightly larger decimal or another proof of gappedness would not meet that
standard. Symmetry-reduced or recursively certified local gaps are possible
approaches, not established improvements. No evidence in this audit yet
shows that the available margin is large enough; benchmark reproduction
and a short numerical test must precede extensive formalization.

\paragraph{Excluded as immediate targets.}
The generic two-block uncle $N^{-2}$ question in the local project notes
is not, on the evidence inspected here, an explicitly located conjecture
from the cited paper. The notes already establish inverse-square bounds
for their scalar model. General two-dimensional PEPS gaps and general
hierarchy completeness require new analysis beyond the existing
one-dimensional formal interfaces. The five-site Heisenberg result in
P3 concerns ground-energy lower bounds, not spectral gaps.

\begin{thebibliography}{9}
\bibitem{LTI}
K.~S.~Rai, I.~Kull, P.~Emonts, J.~Tura, N.~Schuch, and F.~Baccari,
\emph{A Hierarchy of Spectral Gap Certificates for Frustration-Free Spin Systems},
Quantum \textbf{10}, 2065 (2026).
\href{https://doi.org/10.22331/q-2026-04-13-2065}{doi:10.22331/q-2026-04-13-2065}.
Primary text: \url{https://arxiv.org/html/2411.03680v3}.
\bibitem{RMS}
M.~\'A.~Rozm\'an, A.~Moln\'ar, and N.~Schuch,
\emph{Lower Bounds on Spectral Gaps of Parent Hamiltonians via Tensor Networks},
arXiv:2607.19078v1 (2026).
\url{https://arxiv.org/html/2607.19078v1}.
\bibitem{Rao}
S.~Rao, \emph{A convergent hierarchy of spectral gap certificates for qubit
Hamiltonians}, arXiv:2510.08427 (2025).
\url{https://arxiv.org/abs/2510.08427}.
\bibitem{Xu}
X.~Xu, M.~Sch\"otz, J.~Wang, V.~Magron, I.~Klep, O.~Fawzi, and
M.-O.~Renou, \emph{The bulk spectral gap is semi-decidable: a convergent
family of certified upper bounds}, arXiv:2606.03836v2 (2026).
\url{https://arxiv.org/abs/2606.03836v2}.
\bibitem{LM}
M.~Lemm and E.~Mozgunov, \emph{Spectral gaps of frustration-free spin systems
with boundary}, J. Math. Phys. \textbf{60}, 051901 (2019).
\url{https://arxiv.org/abs/1801.08915}.
\bibitem{WY}
S.~Warzel and A.~Young, \emph{The spectral gap of a fractional quantum
Hall system on a thin torus}, J. Math. Phys. \textbf{63},
041901 (2022). \url{https://arxiv.org/abs/2112.13764}.
\end{thebibliography}
\end{document}